\section{Introduction}

Cross-subject EEG decoders are compared across papers whose evaluation protocols use the target subject differently. Some use none of it, some use its unlabelled signals to align feature distributions, some fine-tune on a labelled calibration block. The labels attached to these protocols---zero-calibration, unsupervised transfer, subject-adaptive---are coarse, and the accuracy differences between them are routinely read as the value of the extra access.

That reading requires the comparison to hold everything else constant, and published comparisons rarely can. Adding a calibration block lengthens an epoch, so the adapted model also gets more optimizer steps; the target's share of a batch is whatever the cohort size implies; and if the block supplies the early-stopping signal, model selection has changed too. None of this is visible in a results table, and all of it moves with the effect being claimed. Our own earlier work tracked those quantities rather than the procedure; Section~\ref{sec:controls} reports what replaying it shows.

This paper asks a narrower question that can be answered cleanly. \textbf{On a fixed target test set, what does each way of using the target subject's available trials buy, once the source partition, the per-epoch sample budget, the early-stopping rule and the target sampling rate are held constant?} Each held-out subject's trials are divided once into an available half and a reserved test half, and every condition is scored on the same reserved half. The available half is used in one of four ways: not at all, by Euclidean Alignment, by supervised training, or by both. This crosses two \emph{procedures}, not two independent kinds of access---supervised training reads the available trials' signals as well as their labels---so the grid separates what each procedure adds on top of the other, not the value of signals against labels.

\textbf{Related work.} Partitioning and reporting choices are known to move BCI comparisons \cite{lotte2018review,roy2019deep}, and MOABB standardizes the interface but not the accounting \cite{jayaram2018moabb}: two pipelines can share a loader and a split while giving the target subject different influence over training, normalization and checkpoint choice. Euclidean Alignment reads the evaluated subject's trials but none of their labels and is routinely called zero-calibration \cite{he2020transfer}, while calibration-block fine-tuning consumes labels outright; compact convolutional references \cite{lawhern2018eegnet,schirrmeister2017deep} carry both. Leakage drives irreproducible findings \cite{kapoor2023leakage} and small-sample neuroimaging yields error bars wider than reported \cite{varoquaux2018crossval}, but that work is diagnostic, whereas choosing between protocols needs what a procedure is worth once everything else is matched.

A second question concerns how the result is summarized. A cross-subject paper reports a spread as well as a mean, and the spread is read as between-subject heterogeneity. That holds only if the aggregation unit is the subject. We test it with a control that retrains nothing: the same per-trial predictions are aggregated by true subject and by random groups of the same sizes, so any difference is attributable to the grouping alone.

\textbf{Contributions.} (i) A controlled $2\times2$ of alignment against target supervision on three motor-imagery datasets spanning 9 to 105 subjects, with the source partition, the epoch and the scored trials pinned and checked against the realized indices, and with the exposure quantities that \emph{cannot} simultaneously be pinned reported instead. (ii) A fixed-prediction regrouping benchmark, and a matching variance estimate, for how much of a reported between-subject spread the aggregation unit accounts for. (iii) An audit of what an uncontrolled version of the same comparison reports instead. Code, logs and per-trial predictions: \url{https://github.com/yo666666yo/target-data-access-eeg}. The split, the resampling and alignment itself are standard; what is new is the set of controls under which the comparison is made, and the account of which cannot hold at once.
